\documentclass[11pt]{article}
\usepackage[margin=1in]{geometry}
\usepackage{amsmath,amssymb,amsthm,mathtools,booktabs,longtable,array}
\usepackage[T1]{fontenc}
\usepackage{lmodern,microtype}
\usepackage[hidelinks]{hyperref}
\usepackage{enumitem}
\setlist{nosep}
\newtheorem{theorem}{Theorem}[section]
\newtheorem{lemma}[theorem]{Lemma}
\newtheorem{proposition}[theorem]{Proposition}
\newtheorem{corollary}[theorem]{Corollary}
\theoremstyle{remark}\newtheorem{remark}[theorem]{Remark}
\newcommand{\E}{\mathbb E}
\newcommand{\R}{\mathbb R}
\newcommand{\one}{\mathbf 1}
\newcommand{\tht}{\vartheta}
\newcommand{\sgn}{\sigma}
\DeclareMathOperator{\arsinh}{arsinh}
\title{Random circulant theta: an all-time fixed point and a root-moment criterion}
\author{Sidney Holden\\{\small Center for Computational Biology, Flatiron Institute}\\{\small Simons Foundation}}\date{13 September 2026}
\hypersetup{pdfauthor={Sidney Holden}}
\begin{document}
\maketitle
\begin{abstract}
The random-circulant expectation problem MD-02 asks whether
$\E\tht(G_n)/\sqrt n\to1$ for the inversion-class Bernoulli ensemble, including
even orders. We give an exact fixed-point representation of the complementary
weighted logarithmic barriers. It is valid at every finite barrier time and
comes with a deterministic global convergence bound for Picard iteration.
Complement symmetry converts the target into an explicit moment condition on
one coordinate of the converged fixed point. For the first two Picard iterates
we prove uniform all-orders bounds $n\E(u^{(1)}_0)^2\le2$ and
$n\E(u^{(2)}_0)^2\le68$, without an assumption on the barrier time. These are
finite-iterate theorems, not bounds on the converged root. The required
probabilistic control as the iteration count increases is not proved here.
Consequently this manuscript does not prove or disprove MD-02.
\end{abstract}
\tableofcontents
\newpage

\section{Target and scope}
For $n\ge2$, let $S=-S\subset\mathbb Z_n\setminus\{0\}$ be obtained by
including each inversion class $\{s,-s\}$ independently with probability
$1/2$. The graph $G_S$ has an edge between distinct $i,j$ exactly when
$i-j\in S$. The class $\{n/2\}$ is sampled once when $n$ is even.
The displayed target in MD-02 is
\begin{equation}\label{eq:target}
 \frac{\E\tht(G_S)}{\sqrt n}\longrightarrow1
 \quad\text{through all integer orders.}
\end{equation}
The problem and the repository's status conventions are recorded in
\cite{repo,contributing}. The primary random-circulant theta paper is
\cite{paper}. The earlier supplied manuscripts develop central-path skew
identities, initial moments, and early-time estimates \cite{prior1,prior2}.
Here we use a different representation of the same complementary barriers.
All mathematical results used in the new reduction are proved below.

The scope distinction is essential. The fixed point exists and is uniquely
specified for every deterministic mask, and its first two iterates have
small second moments in the random ensemble. Neither fact proves that the
converged coordinate has a small second moment. No passage from two
iterations to an unbounded number of iterations is made implicitly.

\section{The Fourier probability program}
Let $F$ be the unitary discrete Fourier transform, with
$F_{sj}=n^{-1/2}\exp(-2\pi i sj/n)$. We use the unnormalized transform
$\widehat p=\sqrt n Fp$ for probabilities. Averaging the standard theta
semidefinite program over cyclic translations gives
\begin{equation}\label{eq:lp}
 \tht(G_S)=n\max\{p_0:p\in\R^n_{\ge0},\ 
 \one^Tp=1,\ \widehat p_s=0\ (s\in S)\}.
\end{equation}
For completeness, a translation-invariant primal matrix has entries
$X_{ij}=f(i-j)/n$, with $f(0)=1$ and $f(s)=0$ on edges. Its positive
semidefiniteness is equivalent to nonnegativity of the Fourier eigenvalues.
Writing $p=\widehat f/n$ gives a probability vector and objective
$\sum_s f(s)=np_0$. Real symmetric matrices correspond to even $p$.
Conversely, any real feasible probability can be averaged with its reflection;
this preserves $p_0$ and the constraints. Thus allowing all real probabilities
in \eqref{eq:lp} does not change its value. Complex equalities in that formula
mean both their real and imaginary parts.

Encode the mask by $\sgn_0=0$, $\sgn_s=-1$ for edges and $\sgn_s=+1$
for nonedges, and define
\begin{equation}\label{eq:R}
 R=F^*\operatorname{diag}(\sgn)F,\qquad
 P_0=\frac1n\one\one^T.
\end{equation}
Then $R$ is a real symmetric circulant matrix and
\begin{equation}\label{eq:Ridentities}
 R\one=0,\qquad R^2=I-P_0,\qquad \|R\|_{2\to2}=1.
\end{equation}
The complementary graph changes $R$ to $-R$. Let $r=Re_0$ denote the first
column. Because $R$ is symmetric and circulant, $R_{0j}=r_j$.

\begin{lemma}[Complementary feasible probabilities]\label{lem:orthogonal}
If $p$ is feasible in \eqref{eq:lp} and $q$ is feasible for the complementary
graph, then $p^Tq=1/n$. In particular $np_0\le1/q_0$ when $q_0>0$.
\end{lemma}
\begin{proof}
The nonzero Fourier supports of $p$ and $q$ are disjoint. Parseval's identity
therefore gives $p^Tq=\widehat p_0\widehat q_0/n=1/n$. Nonnegativity yields
$p_0q_0\le1/n$.
\end{proof}

\section{Complementary weighted barriers}
Fix $t\ge1$, and set
\begin{equation}\label{eq:parameters}
 w_0=t,\quad w_j=1\ (j\ne0),\quad W=t+n-1,\quad
 v=w/t,\quad V=W/t=1+(n-1)/t.
\end{equation}
Maximize $\sum_j w_j\log p_j$ over the strictly positive feasible
probabilities for $G_S$. The uniform probability is feasible. The objective
is strictly concave and tends to minus infinity at the boundary, so it has
a unique strictly positive maximizer $p=p(t)$.

\begin{lemma}[Weighted complementary pair]\label{lem:pair}
There is a unique positive complementary probability $q=q(t)$ satisfying
\begin{equation}\label{eq:products}
 p_jq_j=\frac{w_j}{nW}\quad(0\le j<n).
\end{equation}
Moreover $q$ is the weighted-barrier maximizer for the complementary graph.
\end{lemma}
\begin{proof}
The first-order optimality condition says that $w/p$, with division
coordinatewise, belongs to the span of $\one$ and the real Fourier constraint
vectors on $S$. Write $w/p=\lambda\one+h$, with $h$ in the edge Fourier
space. Taking the inner product with $p$ gives $\lambda=W$.
Define $q=(w/p)/(nW)$. It is positive, sums to one, and its Fourier support
is contained in $S\cup\{0\}$. Thus it is complementary feasible and satisfies
\eqref{eq:products}. Conversely $w/q=nWp$ gives its own first-order
optimality condition. Strict concavity proves uniqueness on both sides.
\end{proof}

\begin{proposition}[Deterministic theta interval]\label{prop:theta_interval}
The weighted pair satisfies
\begin{equation}\label{eq:interval}
 np_0\ \le\ \tht(G_S)\ \le\ \frac1{q_0}
       =\frac Wt\,np_0.
\end{equation}
For every circulant graph, $\tht(G_S)\tht(\overline{G_S})=n$.
\end{proposition}
\begin{proof}
The first inequality is primal feasibility and the second follows from
Lemma~\ref{lem:orthogonal}; the equality uses \eqref{eq:products}.
For any optimal complementary probabilities the same lemma gives the product
upper bound $n$. The weighted feasible pair gives the product lower bound
$nt/W$ for every $t$. Letting $t\to\infty$ proves equality.
\end{proof}

\section{The all-time fixed point}
All square roots in what follows are coordinatewise. Define
\begin{equation}\label{eq:map}
 \Phi_{R,t}(u)=R\sqrt{v+u^2}.
\end{equation}

\begin{theorem}[Exact fixed-point representation]\label{thm:fixed}
For every mask and every finite $t\ge1$, the equation
\begin{equation}\label{eq:fixed}
 u=R\sqrt{v+u^2}
\end{equation}
has exactly one solution $u^*=u^*(R,t)$ in $\R^n$. Put
$c=\sqrt{v+(u^*)^2}$. The complementary barrier pair is
\begin{equation}\label{eq:reconstruct}
 p=\frac{c+u^*}{\sqrt{nV}},\qquad
 q=\frac{c-u^*}{\sqrt{nV}}.
\end{equation}
In particular,
\begin{equation}\label{eq:rootformula}
 \sqrt{\frac tW}\bigl(\sqrt{1+(u^*_0)^2}+u^*_0\bigr)
 \le\frac{\tht(G_S)}{\sqrt n}
 \le\sqrt{\frac Wt}\bigl(\sqrt{1+(u^*_0)^2}+u^*_0\bigr).
\end{equation}
Complementation changes $u^*$ to $-u^*$.
\end{theorem}
\begin{proof}
Start with Lemma~\ref{lem:pair} and define
$c=\sqrt{nV}(p+q)/2$, $u=\sqrt{nV}(p-q)/2$. Equation
\eqref{eq:products} gives $c_j^2-u_j^2=v_j$, and $c_j>0$.
At nonzero Fourier frequencies the two probabilities have complementary
supports, so $u=Rc$. This proves existence.

Conversely suppose \eqref{eq:fixed} holds, and put $c=\sqrt{v+u^2}$.
By \eqref{eq:Ridentities}, $\one^Tu=0$ and
\[
 \|u\|^2=\|c\|^2-\frac{(\one^Tc)^2}{n}.
\]
Since $\|c\|^2-\|u\|^2=V$ and $c$ is positive,
$\one^Tc=\sqrt{nV}$. Thus \eqref{eq:reconstruct} defines positive
probabilities. Their Fourier supports are complementary, and their
coordinate products are $v_j/(nV)=w_j/(nW)$. They satisfy the
first-order conditions of Lemma~\ref{lem:pair}, so uniqueness of the barrier
maximizer proves uniqueness of $u$.

At coordinate zero, $c_0=\sqrt{1+u_0^2}$ and
$np_0/\sqrt n=\sqrt{t/W}(c_0+u_0)$. Proposition~\ref{prop:theta_interval}
gives \eqref{eq:rootformula}. Finally the square-root map is even: if $u$ is
fixed for $R$, then $-u$ is fixed for $-R$. Uniqueness proves the last claim.
\end{proof}

The barrier logarithm is therefore exactly
\begin{equation}\label{eq:logroot}
 L_t=\arsinh(u^*_0),\qquad e^{L_t}=\sqrt{1+(u^*_0)^2}+u^*_0.
\end{equation}
This connects the fixed point to the logarithmic coordinate used in the prior
central-path manuscripts. It does not estimate that coordinate.

\section{Global convergence with an explicit error bound}
Let $u^{(0)}=0$ and $u^{(k+1)}=\Phi_{R,t}(u^{(k)})$.

\begin{theorem}[Deterministic convergence]\label{thm:convergence}
For every mask, every $t\ge1$, and every $k\ge0$,
\begin{equation}\label{eq:convergence}
 \|u^{(k)}-u^*\|_2\le Dq^k,\qquad
 D=\sqrt{\frac{nW}{2t}},\quad
 q=\sqrt{\frac{2nW}{1+2nW}}<1.
\end{equation}
Consequently Picard iteration converges globally from zero. These bounds are
exact-arithmetic mathematical bounds; ordinary floating-point computation
needs additional roundoff analysis to turn them into interval certificates.
\end{theorem}
\begin{proof}
For a zero-sum vector $u$, its positive and negative parts have equal total
mass. Hence $\|u\|_2^2\le\|u\|_1^2/2$. At the fixed point,
$\|u^*\|_1\le\one^T\sqrt{v+(u^*)^2}=\sqrt{nV}$, which proves
$\|u^*\|_2\le D$.

Every scalar map $x\mapsto\sqrt{v_j+x^2}$ is one-Lipschitz. Together with
$\|R\|\le1$, this shows that $\Phi$ is nonexpansive, and induction gives
$\|u^{(k)}-u^*\|\le D$. Thus all iterates have norm at most $2D$.
On the interval $[-2D,2D]$ the scalar derivatives have absolute value at most
\[
 \frac{2D}{\sqrt{t^{-1}+4D^2}}=q.
\]
The line segment joining an iterate and the fixed point stays within these
coordinate intervals. The mean-value inequality therefore improves
nonexpansiveness to contraction by $q$ for these pairs, proving
\eqref{eq:convergence} by induction.
\end{proof}

\begin{corollary}[A finite-iteration version of the missing estimate]\label{cor:finite}
Let $t=n^2$ and
\[
 k_n=\left\lceil4(1+2n(t+n-1))\log n\right\rceil.
\]
Then $\|u^{(k_n)}-u^*\|_2=O(n^{-3/2})$, uniformly in the mask.
A proof that $\E(u^{(k_n)}_0)^2=o(1)$ would therefore establish the sufficient
root-moment condition in Theorem~\ref{thm:criterion} below.
\end{corollary}
\begin{proof}
Writing $A=2nW$, one has
$\log q=-\tfrac12\log(1+1/A)\le-1/(2(1+A))$.
Thus $q^{k_n}\le n^{-2}$, while $D=O(\sqrt n)$.
The last assertion follows from the triangle inequality in $L^2$.
\end{proof}

The iteration count in this corollary is of order $n^3\log n$. Proving a
moment estimate for two iterations is not a substitute for proving it at this
count. The convergence theorem itself is entirely deterministic and supplies
no such probabilistic estimate.

\section{An exact expectation criterion at the endpoint}
The random mask has the same distribution as its complement. All fixed
points in a fixed dimension are bounded by Theorem~\ref{thm:convergence};
in particular the expectations below exist.

\begin{theorem}[Root-moment criterion]\label{thm:criterion}
For every $n\ge2$ and $t\ge1$, set
\[
 B_{n,t}=\E\bigl[\sqrt{1+(u^*_0(R,t))^2}-1\bigr]\ge0.
\]
Then
\begin{equation}\label{eq:exactcriterion}
 \sqrt{\frac tW}(1+B_{n,t})
 \le\frac{\E\tht(G_S)}{\sqrt n}
 \le\sqrt{\frac Wt}(1+B_{n,t}).
\end{equation}
For any sequence $t_n/n\to\infty$, the MD-02 limit holds if and only if
$B_{n,t_n}\to0$. Equivalently it holds if and only if
\begin{equation}\label{eq:necessarysufficient}
 \E\min\{(u^*_0(R,t_n))^2,\ |u^*_0(R,t_n)|\}\longrightarrow0.
\end{equation}
The simpler condition $\E(u^*_0(R,t_n))^2\to0$ is sufficient.
\end{theorem}
\begin{proof}
Complementation negates $u^*$, so $\E u^*_0=0$. Take expectations in
\eqref{eq:rootformula} to obtain \eqref{eq:exactcriterion}. If $t_n/n\to\infty$,
both deterministic prefactors tend to one. Since $B\ge0$, the two-sided
inequality proves both directions of the equivalence.

For every real $x$,
\[
 \frac1{1+\sqrt2}\min\{x^2,|x|\}
 \le\sqrt{1+x^2}-1
 \le\min\{x^2,|x|\}.
\]
This follows by rationalizing $\sqrt{1+x^2}-1$ and distinguishing
$|x|\le1$ from $|x|\ge1$. It proves \eqref{eq:necessarysufficient}.
Finally $\sqrt{1+x^2}-1\le x^2/2$ proves sufficiency of the second moment.
\end{proof}

\begin{corollary}[A quantitative sufficient bound]\label{cor:rate}
If $t=n^2$ and $\E(u^*_0)^2\le C/n$, then
\[
 1\le\frac{\E\tht(G_S)}{\sqrt n}
 \le 1+\frac{C+1}{2n}+\frac{C}{4n^2}.
\]
In particular, a uniform bound
$\sup_{k\ge0}n\E(u^{(k)}_0(R,n^2))^2\le C$ would prove MD-02.
Neither bound in this corollary is established in this manuscript.
\end{corollary}
\begin{proof}
For the lower bound, pair a graph with its complement and use
$\tht(G)\tht(\bar G)=n$ and the arithmetic--geometric mean inequality.
For the upper bound apply $\sqrt{1+y}\le1+y/2$ twice in
\eqref{eq:exactcriterion}, using $(n-1)/n^2\le1/n$.
For the final implication, convergence of the iterates and Fatou's lemma
transfer a hypothetical uniform bound to $u^*$.
\end{proof}

\section{Uniform moments of the first two Picard roots}
This section proves the finite-iterate estimates without assuming independence
of the Fourier entries. Let $\mathcal C_n$ be the nonzero inversion classes,
and $m_c=|c|\in\{1,2\}$. Define
\[
 h_c=F^*\operatorname{diag}(1_c)Fe_0,
 \qquad r=\sum_{c\in\mathcal C_n}\sgn_c h_c.
\]
These vectors are real, mutually orthogonal, and satisfy
\begin{equation}\label{eq:hc}
 \|h_c\|_2^2=\frac{m_c}{n},\qquad
 (h_c)_0=\frac{m_c}{n},\qquad \one^Th_c=0.
\end{equation}
If $A$ is the matrix with columns $h_c$, then
\begin{equation}\label{eq:Ac}
 \|A\|_{2\to2}^2\le\frac2n,\qquad
 \|r\|_2^2=\frac{n-1}{n},\qquad
 \sum_c m_c^2\le2(n-1).
\end{equation}

\begin{theorem}[All-time finite-iterate bounds]\label{thm:two}
For $a=1-t^{-1/2}$, every $n\ge2$, and every finite $t\ge1$,
\begin{align}
 n\E(u^{(1)}_0)^2
 &=a^2\begin{cases}2-2/n,&n\text{ odd},\\2-3/n,&n\text{ even},\end{cases}
 \label{eq:first}\!\\[-1ex]
 n\E(u^{(2)}_0)^2&\le68a^2.\label{eq:second}
\end{align}
Both roots have mean zero. The same conclusions hold at the limit $t\to\infty$
for these finite iterates, interpreted by continuity.
\end{theorem}
\begin{proof}
Since $\sqrt v=t^{-1/2}\one+ae_0$ and $R\one=0$,
$u^{(1)}=ar$. Independence of the class signs gives
$\E r_0^2=n^{-2}\sum_c m_c^2$. For odd $n$ the sum is $2(n-1)$; for even
$n$ it is $2(n-2)+1=2n-3$. This proves \eqref{eq:first}.

Put $b=t^{-1}$ and
\[
 f(x)=x\sqrt{b+a^2x^2},\qquad
 g(x)=f(x)-\sqrt b\,x,\qquad H=\sum_{j=0}^{n-1}g(r_j).
\]
Because $\sum_j r_j=0$, also $H=\sum_jf(r_j)$. Differentiation yields
\begin{equation}\label{eq:gder}
 |g'(x)|\le2a|x|,\qquad |g''(x)|\le2a.
\end{equation}
For the second bound, when $a,b>0$ set $z=a|x|/\sqrt b$ and use
$z(3+2z^2)/(1+z^2)^{3/2}\le2$. The cases $a=0$ and the limit $b\downarrow0$
follow directly or by continuity.

Flip one class sign $\sgn_c$. The corresponding kernel changes from $r$
to $r-2\sgn_ch_c$. Taylor's theorem and \eqref{eq:gder} give
\[
 H(\sgn^{(c)})-H(\sgn)
   =-2\sgn_c\langle h_c,g'(r)\rangle+\epsilon_c,
 \qquad |\epsilon_c|\le4a\|h_c\|_2^2=4am_c/n.
\]
Thus, pointwise on the sign cube,
\begin{align*}
 \sum_c\bigl(H(\sgn^{(c)})-H(\sgn)\bigr)^2
 &\le8\|A^Tg'(r)\|_2^2+\frac{32a^2}{n^2}\sum_cm_c^2\\
 &\le\frac{16}{n}\cdot4a^2\|r\|_2^2+\frac{64a^2(n-1)}{n^2}
 \le\frac{128a^2}{n}.
\end{align*}
The Rademacher Poincare inequality gives
\begin{equation}\label{eq:Hvar}
 \operatorname{Var}H\le\frac14\E\sum_c
     (H(\sgn^{(c)})-H(\sgn))^2\le\frac{32a^2}{n}.
\end{equation}
For a self-contained justification of the factor $1/4$, expand a function
of independent signs in the orthonormal Walsh basis. Its variance is the sum
of squares of all nonconstant coefficients. The right side before the last
inequality weights each such coefficient by the size of its nonempty index
set, and is therefore no smaller.

The exceptional root weight gives
\[
 u^{(2)}_0=H+\delta(r_0),\qquad
 \delta(x)=x\left(\sqrt{1+a^2x^2}-\sqrt{b+a^2x^2}\right).
\]
The expression in parentheses lies between zero and $1-\sqrt b=a$, so
$|\delta(x)|\le a|x|$. Both $H$ and $\delta(r_0)$ are odd under global
sign reversal, so have mean zero. Therefore
\[
 \E(u^{(2)}_0)^2\le2\E H^2+2\E\delta(r_0)^2
 \le\frac{64a^2}{n}+\frac{4a^2}{n}=\frac{68a^2}{n}.
\]
Continuity in $b\downarrow0$ and the finite size of the mask space for each
fixed $n$ justify the endpoint statement for these two iterates.
\end{proof}

\begin{corollary}[First-root central limit]\label{cor:clt}
If $t_n\to\infty$, then $\sqrt n\,u^{(1)}_0(R,t_n)$ converges in distribution
to a centered Gaussian of variance $2$, along all integer orders.
\end{corollary}
\begin{proof}
The first root is $a_n\sum_cm_c\sgn_c/\sqrt n$. Its variance tends to $2$
and the largest summand has magnitude at most $2/\sqrt n$. The Lindeberg
condition for this independent triangular array is immediate.
\end{proof}

Theorem~\ref{thm:two} makes no Gaussian approximation and treats the singleton
exactly. The constant $68$ is a valid explicit bound, not a claimed optimal
constant. No central limit theorem for the second or later root is needed
for, or asserted by, this theorem.

\section{Why the two-iterate theorem is not a solution}
There are two distinct limits: $n\to\infty$ and $k\to\infty$.
Theorem~\ref{thm:two} bounds the root when $k=1,2$ for all barrier times.
Theorem~\ref{thm:criterion} concerns $k=\infty$. The deterministic error
bound in Theorem~\ref{thm:convergence} does not bridge this difference at
$k=2$; at $t=n^2$ its bound is of order $\sqrt n$ there.

A finite iterate also cannot be silently inserted into
\eqref{eq:rootformula}. The identity $\one^Tc=\sqrt{nV}$ used to reconstruct
feasible normalized probabilities relies on the fixed-point equation itself.
It need not hold for a finite iterate.

Even the smallest deterministic example distinguishes the second iterate
from the fixed point. For the empty graph on two vertices,
$R=I-P_0$. In the limit $t\to\infty$,
\[
 u^{(1)}_0=\frac12,\qquad
 u^{(2)}_0=\frac{\sqrt5-1}{4},\qquad
 u^*_0=\frac1{2\sqrt2}.
\]
The middle and final numbers are unequal. The first two formulas follow
directly from the map. For the final formula, the one-variable fixed-point
equation is $u=(\sqrt{1+u^2}-u)/2$, which has the displayed positive
solution. This is an illustration of the logical distinction, not a
counterexample to the random asymptotic statement.

Likewise a bad deterministic mask at large $n$ cannot by itself refute
MD-02: its probability and its contribution to the expectation have to be
estimated. The empty and complete masks each have probability
$2^{-\lfloor n/2\rfloor}$. Their combined contribution to the normalized
expectation is at most $(n+1)2^{-\lfloor n/2\rfloor}/\sqrt n$, which tends
to zero.

The exact outstanding task in this formulation is
\begin{equation}\label{eq:missing}
 \E\bigl[\sqrt{1+(u^*_0(R,n^2))^2}-1\bigr]\longrightarrow0.
\end{equation}
A bound $\E(u^*_0)^2=O(1/n)$ would prove substantially more than necessary,
but it is not established by the two-iterate argument. The evolving
nonlinearity depends on the same random mask as $R$; later applications
cannot be treated as multiplication by an independent random operator.

\section{Reproducibility and claim audit}
The accompanying implementation applies $R$ by FFT, computes finite Picard
iterates, reports the deterministic exact-arithmetic error formula, and
numerically solves the fixed point using a continuation in $t$. A separate
Fourier linear program checks the theta interval. Root-solver residuals are
recorded explicitly; these solves are not interval-arithmetic certificates.

The validation script enumerates inversion-class masks at orders $2$ through
$12$ for first- and second-root moments, checks weighted reconstruction and
complement symmetry against Fourier linear programs on smaller instances,
and records Monte Carlo diagnostics separately. The actual run outcome is
in \texttt{results/validation\_summary.json}. An implemented test or a
successful finite experiment is not a proof of \eqref{eq:missing}.

\begin{longtable}{p{0.32\textwidth}p{0.18\textwidth}p{0.41\textwidth}}
\toprule
Claim & Status here & Scope\\\midrule
Exact all-time fixed point & Proved & Every mask, every finite $t\ge1$; Theorem~\ref{thm:fixed}.\\
Global Picard convergence & Proved & Explicit deterministic error bound; Theorem~\ref{thm:convergence}.\\
Expectation criterion & Proved equivalence & Concerns the converged root; Theorem~\ref{thm:criterion}.\\
First two root moments & Proved & All integer orders and all barrier times; Theorem~\ref{thm:two}.\\
First-root Gaussian limit & Proved & Independent triangular-array limit; Corollary~\ref{cor:clt}.\\
Converged-root moment & Not proved & Equation~\eqref{eq:missing} is the remaining task.\\
MD-02 expectation limit & Not proved or disproved & No substantive asymptotic subcase of the displayed target is established.\\
Repository status recommendation & OPEN & Supporting results alone do not establish a displayed-target case.\\
\bottomrule
\end{longtable}

The mathematical proofs are not accompanied by a claim of independent peer
review or historical priority. Earlier proof archives, when available, are
preserved as provenance rather than incorporated as unexamined hypotheses.
The new theorems above are self-contained.

\paragraph{Submission preparation.} AI assistance was used for submission preparation and a separate informal Codex AI-agent audit. This is not external human peer review or formal verification. No Lean verification was performed.

\begin{thebibliography}{9}
\bibitem{repo}
OpenProblemsInNLA. MD-02, random circulant Lovasz theta, problem statement.
\url{https://github.com/ajt60gaibb/OpenProblemsInNLA/tree/main/matrix-discrepancy-and-optimization/MD-02}.
\bibitem{contributing}
OpenProblemsInNLA. Repository README and contribution guidelines.
\url{https://github.com/ajt60gaibb/OpenProblemsInNLA};
\url{https://github.com/ajt60gaibb/OpenProblemsInNLA/blob/main/CONTRIBUTING.md}.
\bibitem{paper}
A.~S. Bandeira, J. B{\l}asiok, D. Dmitriev, U. Faure, A. Kireeva, and
D. Kunisky. \emph{The Lov\'asz number of random circulant graphs}.
arXiv:2502.16227, 2025. \url{https://arxiv.org/abs/2502.16227}.
\bibitem{prior1}
Supplied continuation archive, \texttt{MD02\_continuation\_results.zip}:
\emph{Central-path continuation, sharp initial moments, and certified
obstructions}. Sections 2--4 give the Fourier formulation and weighted
complementary barriers; subsequent sections give initial-moment results and
obstructions.
\bibitem{prior2}
Supplied continuation archive, \texttt{MD02\_extended\_proofs.zip}:
\emph{Random circulant theta: initial limit laws and early central-path
control}. Sections 2 and 8 give barrier notation and the full-scale gap;
Theorem 7.1 gives early-time control.
\end{thebibliography}
\end{document}
